\begin{abstract}
Subword tokenizers learn their vocabulary from frequency statistics, so languages with
less training text pay more tokens for the same content. Classical BPE minimises the total
token count and ignores this disparity, while Parity-aware BPE equalises the cost across
languages at some loss of compression. Methods between these ends are heuristic and
state no objective. To address this, we show that both ends
are special cases of one objective, a power mean of per-language token costs with exponent
$p$. Its greedy step weights each language's pair counts by its current cost raised to the
power $p-1$, giving classical BPE at $p=1$ and Parity-aware BPE as $p\to\infty$. We call the
tokenizers in between $p$-fair BPE and train them exactly by lazy evaluation, about 20 times
faster than full rescoring. On Hindi, Marathi, Gujarati and English, raising $p$
traces a monotone compression--parity front in all five training samples. At $p=8$, the
Gini coefficient of per-language cost falls by 40\% for a 0.6\% loss in compression. The
% TODO (E3 in experiments.md): this "no language worse off" statement holds for the four-language
% setting only. In the 14-language pilot (reference_results/l14_preview) p-fair raises the cheapest
% languages (Bengali, Hindi) like every fairness rule does; what stays true is that p-fair never
% raises a language that costs MORE than the classical mean, whereas the hybrid does (Gujarati,
% Malayalam). Reword to the under-served form once E3 has error bars (paper_tables.py underserved).
front matches the hybrid baseline, but unlike the hybrid, $p$-fair BPE never raises an
Indian language's cost above its classical cost. Finally, we show that the cost unit and
data balance decide the size and even the direction of the disparity.
\end{abstract}